\sgaStableTarget{sga6.e06.sec.2.3}2.3.
We shall need to apply the functor @@M0@@ to arbitrary graded @@M1@@-Modules, and not only to quasi-coherent graded @@M2@@-Modules as in the construction of EGA II 3.2. We shall therefore indicate briefly how this construction can be carried out in the general case.

\sgaStableTarget{sga6.e06.para.2-3-1}2.3.1. First let @@M3@@ be a scheme, @@M4@@ a quasi-coherent @@M5@@-algebra, @@M6@@, and @@M7@@ the structural morphism. From the isomorphism:
@@M8@@
we deduce a canonical homomorphism:
\sgaStableTarget{sga6.e06.eq.2-3-1-1}@@M9@@
where @@M10@@ denotes inverse image in the sense of sheaves of sets; in other words, we have a morphism of ringed spaces:
\sgaStableTarget{sga6.e06.eq.2-3-1-2}@@M11@@
From it we deduce an ``inverse-image'' functor @@M12@@ from the category of @@M13@@-Modules to the category of @@M14@@-Modules, left adjoint to the ``direct-image'' functor @@M15@@ from the category of @@M16@@-Modules to the category of @@M17@@-Modules. We therefore have:
\sgaStableTarget{sga6.e06.eq.2-3-1-3}@@M18@@
for every @@M19@@-Module @@M20@@.

Since the homomorphism (2.3.1.1) is flat (since on each stalk it corresponds to a localization), the functor @@M21@@ thus defined is an exact functor (cf [4], p. 164).

If @@M22@@ is a quasi-coherent @@M23@@-Module, this construction gives the Module @@M24@@ defined in EGA II 1.4.3. Indeed, the isomorphism:
